% ============================================================
\section{无限集：偏好仍可能被效用表示}
% ============================================================

有限集证明的两条腿——“取最差元”与“有限计数”——在无限集上都失效。
但\textbf{无限本身不构成障碍}：$\R_+$ 上的通常大小关系由 $u(x)=x$ 平凡表示。
真正的问题是：哪些无限结构允许数值化？本节按基数从小到大推进。

\subsection{可数无限集：权数构造}

\begin{theorem}[可数集表示定理]\label{thm:countable}
设 $X$ 为至多可数集，$\succeq$ 为 $X$ 上的弱序。则 $\succeq$ 可效用表示。
\end{theorem}

\begin{proof}
设 $X=\{x_1,x_2,x_3,\dots\}$ 为一个枚举（有限情形已由定理~\ref{thm:finite} 覆盖）。
给第 $n$ 个元素分配\textbf{几何权数} $w_n=2^{-n}$，定义
\[
u(x)\ \coloneqq\ \sum_{n:\,x\succ x_n} 2^{-n}\ -\ \sum_{n:\,x_n\succ x} 2^{-n},
\]
即“被 $x$ 击败者的权数和”减去“击败 $x$ 者的权数和”。两级数正项且总权 $\sum 2^{-n}\le 1$，
故 $u(x)$ 是良定义的实数。验证表示性。

\textbf{(i) $x\succ y\Rightarrow u(x)>u(y)$}：定义差集指标
\[
S(x,y)\ \coloneqq\ \{n:\ x\succeq x_n\ \text{且}\ y\not\succeq x_n\}\ \cup\ \{n:\ x_n\succeq y\ \text{且}\ x_n\not\succeq x\}.
\]
由传递性，$x\succ y$ 时“$x$ 的下截面”严格包含“$y$ 的下截面”、
“$y$ 的上截面”严格包含“$x$ 的上截面”，故 $u(x)-u(y)$ 是两个正项级数之差，
且每个指标 $n$ 对差值的贡献非负、至少一个指标贡献严格为正（$y$ 自身满足 $x\succ y$，
对应指标贡献 $2^{-n(y)}$ 于第一部分）。于是
\[
u(x)-u(y)\ \ge\ 2^{-n(y)}\ >\ 0.
\]

\textbf{(ii) $u(x)>u(y)\Rightarrow x\succ y$}：由完备性，若 $y\succeq x$，
则 $x\succ y$ 不成立，且分两种情形：$x\sim y$ 时两者的上、下截面完全相同
（$x\sim y$ 与传递性给出 $x\succeq z\iff y\succeq z$），故 $u(x)=u(y)$，矛盾；
$y\succ x$ 时由 (i) 得 $u(y)>u(x)$，矛盾。完备性排除了第三种可能，故 $x\succ y$。

合起来 $x\succeq y\iff u(x)\ge u(y)$。
\end{proof}

\begin{remark}[权数构造的思想]
给每个元素一个“可求和的权重”，把\textbf{集合包含关系}翻译成\textbf{数值大小关系}。
几何权数 $2^{-n}$ 保证任意子集的权和有限，差值的正负完全由包含方向决定。
同一个思想将在 §3.4 升级为：先把一个可数序稠密集“保序编号”进 $(0,1)$，再延拓到整个 $X$。
\end{remark}

\begin{example}[无限阶梯]
$X=\N$，偏好 $1\prec 2\prec 3\prec\cdots$（越大越好）。
权数构造给出 $u(k)=\sum_{n<k}2^{-n}=1-2^{-(k-1)}$，取值于 $[0,1)$；
当然 $u(k)=k$ 也是合法表示——再次印证表示不唯一（命题~\ref{prop:ordinal}），
严格递增变换 $f(t)=-\ln(1-t)$（$t<1$）把前者映为后者的平移。
\end{example}

\subsection{连续统：连续偏好与连续效用}

对 $X=\R^n_+$ 这类不可数消费集，权数构造无法直接套用（不可数正项和发散）。
经济学中的标准出路是\textbf{拓扑条件}：用连续性替代可数性。

\begin{definition}[连续偏好]\label{def:cont-pref}
设 $X\subseteq\R^n$。弱序 $\succeq$ \textbf{连续}，当且仅当对每个 $x\in X$，
上轮廓集 $\{y\in X:y\succeq x\}$ 与下轮廓集 $\{y\in X:x\succeq y\}$ 都是 $X$ 中的闭集。
等价地：若 $y^k\succeq x$（或 $x\succeq y^k$）且 $y^k\to y$，则 $y\succeq x$（或 $x\succeq y$）。
\end{definition}

\begin{remark}[连续性的序含义]
“极限运算不翻转偏好方向”：序列中每一项都至少和 $x$ 一样好，极限点不会突然变得比 $x$ 差。
词典序偏好正是破坏这一点的典型。取
\[
y^k=\Bigl(\frac{1}{k},\,1\Bigr),\qquad x=(0,\,2).
\]
对每个 $k$，$y^k$ 的第一坐标 $1/k>0$，故 $y^k\succ_L x$，即 $y^k$ 全落在 $x$ 的上轮廓集
$U(x)=\{z:z\succeq_L x\}$ 中；但 $y^k\to(0,1)$，而 $(0,1)\prec_L(0,2)=x$，极限点\textbf{掉出}了 $U(x)$。
故 $U(x)$ 不是闭集，$\succeq_L$ 不连续。
这从拓扑角度重新解释了定理~\ref{thm:lex}：不连续的偏好不可能有连续效用表示；
而 §\ref{subsec:lex} 的基数论证更强——它连不连续的效用表示都没有。
\end{remark}

\begin{theorem}[连续偏好的效用表示定理]\label{thm:continuous-rep}
设 $X\subseteq\R^n$ 为非空闭凸集（或更一般地，连通可分拓扑空间），
$\succeq$ 为 $X$ 上连续的弱序。则存在\textbf{连续}效用函数 $u:X\to\R$ 表示 $\succeq$。
\end{theorem}

证明分三步，核心步骤（连续编号）在 §3.4 完整展开：

\begin{enumerate}
  \item \textbf{可数序稠密子集的存在性}：$X$ 可分（含可数稠密集 $D$），
        连续性把“$D$ 稠密”升级为“$D$ 序稠密”：
        若 $x\succ y$，连续性给出开球 $B(x,\varepsilon)$ 内所有点仍 $\succ y$，取 $d\in D\cap B(x,\varepsilon)$，
        则 $x\succeq d\succeq y$（严格说需多一步插值，见 §3.4 引理~\ref{lem:interp}）。
  \item \textbf{连续编号}：把 $D$ 保序地映成 $(0,1)$ 的稠密子集（§3.4 定义~\ref{def:cont-num}）。
  \item \textbf{延拓与连续性}：用 $u(x)=\sup\{g(d):d\in D,\ x\succeq d\}$ 延拓到 $X$；
        偏好连续 $\Rightarrow$ $u$ 的上下轮廓集闭 $\Rightarrow$ $u$ 连续（§3.4 定理~\ref{thm:u-cont}）。
\end{enumerate}

\begin{remark}[与必要条件定理的关系]
定理~\ref{thm:dense-necessary} 说可数序稠密子集\textbf{必要}；
本定理的证明路线说明在拓扑配合下它\textbf{也充分}。
事实上去掉全部拓扑假设后，“弱序 + 可数序稠密子集”本身就是充要条件（§3.4 完成其充分性证明）——
连续性是“让 $X$ 自带一个现成的可数序稠密集”的经济学方便假设，
并且额外买回了效用函数的连续性（这对最优化理论至关重要：
连续效用 + 紧预算集 $\Rightarrow$ 最优解存在）。
\end{remark}

\subsection{单调性给出的显式表示 \texorpdfstring{$\star$}{*}}

对单调偏好，表示可以更具体。

\begin{proposition}[单调连续偏好的对角线构造]\label{prop:diagonal}
设 $\succeq$ 为 $\R^n_+$ 上连续、强单调的弱序。
则每个无差异类与对角线 $\{t\mathbf{1}:t\ge0\}$（$\mathbf{1}=(1,\dots,1)$）恰交于一点，
从而
\[
u(x)\ \coloneqq\ t\ \text{使得}\ x\sim t\mathbf{1}
\]
是表示 $\succeq$ 的连续效用函数。
\end{proposition}

\begin{proof}[证明要点]
\textbf{存在性}：设 $\bar t=\max_i x_i$。由强单调性 $\bar t\mathbf{1}\succeq x\succeq 0\cdot\mathbf{1}$。
考察 $g(t)=\mathbf{1}_{t\mathbf{1}\succeq x}$：连续性使 $\{t:t\mathbf{1}\succeq x\}$ 为闭区间
（含 $\bar t$，下确界记为 $t^*$），$\{t:x\succeq t\mathbf{1}\}$ 亦为闭区间（含 $0$，上确界记为 $t^{*}$）。
完备性迫使两闭区间覆盖 $[0,\bar t]$，交点处 $t\mathbf{1}\sim x$。
\textbf{唯一性}：$t_1<t_2\Rightarrow t_2\mathbf{1}\gg t_1\mathbf{1}\Rightarrow t_2\mathbf{1}\succ t_1\mathbf{1}$（强单调），
故无差异类至多交对角线一次。
\textbf{表示性}：$x\succeq y\iff u(x)\mathbf{1}\sim x\succeq y\sim u(y)\mathbf{1}\iff u(x)\mathbf{1}\succeq u(y)\mathbf{1}\iff u(x)\ge u(y)$，
末步用强单调性。
\textbf{连续性}：$u$ 的上、下轮廓集即 $\succeq$ 的上、下轮廓集，由连续性为闭集。
\end{proof}

\begin{center}
\begin{tikzpicture}[scale=1.15]
  \draw[axisline] (0,0) -- (5.2,0) node[below left, axislabel] {$x_1$};
  \draw[axisline] (0,0) -- (0,5.2) node[below left, axislabel] {$x_2$};
  % diagonal
  \draw[thick, dashed, AxisBlue] (0,0) -- (4.9,4.9) node[right, align=left, axislabel] {对角线\\ $t\mathbf{1}$};
  % indifference curves
  \draw[thick, MidBlue, domain=0.55:4.6, smooth, variable=\x] plot ({\x},{1.0/\x*2.4});
  \draw[thick, MidBlue, domain=0.85:4.6, smooth, variable=\x] plot ({\x},{1.0/\x*4.2});
  \draw[thick, MidBlue, domain=1.3:4.6, smooth, variable=\x] plot ({\x},{1.0/\x*6.4});
  % intersection points
  \node[nodeblue, label=below right:{$u(x')\mathbf{1}$}] at (1.549,1.549) {};
  \node[nodeblue, label=below right:{$u(x)\mathbf{1}$}] at (2.049,2.049) {};
  \node[nodeblue] at (2.53,2.53) {};
  % a bundle x
  \node[nodeblue, label=right:{$x$}] at (3.6,1.1667) {};
  \node[align=left] at (4.3,4.2) {无差异曲线与对角线\\ 恰交于一点（强单调）};
\end{tikzpicture}
\end{center}
\figexplan{对角线 $t\mathbf{1}$ 是构造效用值的标尺：强单调性保证每条无差异曲线与对角线恰交于一点，该交点 $u(x)\mathbf{1}$ 的数值即定义为 $x$ 的效用。图中三条双曲型无差异曲线由低到高排列，越靠右上效用越高；$x$ 所在无差异曲线与对角线的交点给出 $u(x)$。}

命题~\ref{prop:diagonal} 的构造把“效用值”定义为\textbf{等值的对角线消费束的坐标}，
这是福利分析中“等值变化”的雏形。同一思想在 §3.7 拟线性情形将精确化为
$u(t,y)=t+v(y)$，其中 $v(y)$ 就是“用计价物度量的 $y$ 的价值”。
